% Part: lambda-calculus
% Chapter: church-rosser
% Section: parallel-beta-eta-reduction

\documentclass[../../../include/open-logic-section]{subfiles}

\begin{document}

\olfileid{lam}{cr}{pbe}

\olsection{هممهاله $\beta\eta$-راکمول}

په دې برخه کښې د هممهاله $\beta\eta$-راکمول د چرچ--روسر خاصيت
ثابتوو۔ دا د $\beta\eta$-راکمول اړوند هممهاله راکمول دے۔

\begin{defn}[هممهاله $\beta\eta$-راکمول، $\beredpar$] \ollabel{defn:beredpar}
  پر ترمونو \emph{هممهاله $\beta\eta$-راکمول} ($\beredpar$)
  په استقرايي ډول داسې تعريفېږي:
  \begin{enumerate}
    \item \ollabel{defn:beredpar1} $x \beredpar x$.
    \item \ollabel{defn:beredpar2} که $N \beredpar N'$ وي، نو $\lambd[x][N] \beredpar
      \lambd[x][N']$.
    \item \ollabel{defn:beredpar3} که $P \beredpar P'$ او $Q \beredpar Q'$ وي، نو $PQ \beredpar
      P'Q'$.
    \item \ollabel{defn:beredpar4} که $N \beredpar N'$ او $Q \beredpar Q'$ وي، نو
      $(\lambd[x][N])Q \beredpar \Subst{N'}{Q'}{x}$.
    \item \ollabel{defn:beredpar5} که $N \beredpar N'$ وي، نو $\lambd[x][Nx]
      \beredpar N'$، په دې شرط چې $x \notin FV(N)$ وي۔
  \end{enumerate}
\end{defn}

\begin{thm}\ollabel{thm:refl}
  $M \beredpar M$.
\end{thm}

\begin{proof}
  تمرين۔
\end{proof}

\begin{prob}
  \olref[lam][cr][pbe]{thm:refl} ثابت کړئ۔
\end{prob}

\begin{defn}[$\beta\eta$-بشپړ راکمول]\ollabel{defn:becd}
  د~$M$ \emph{$\beta\eta$-بشپړ راکمول}~$\becd{M}$ داسې
  تعريفېږي:
  \begin{align}
    \becd{x} &= x \ollabel{defn:becd1} \\
    \becd{(\lambd[x][N])} &= \lambd[x][\becd{N}] \ollabel{defn:becd2}\\
    \becd{(PQ)} &= \becd{P}\becd{Q} && \text{که $P$ د $\lambd$ تجريد نۀ وي} 
    \ollabel{defn:becd3} \\
    \becd{((\lambd[x][N])Q)} &= \Subst{\becd{N}}{\becd{Q}}{x}
                             \ollabel{defn:becd4} \\
    \becd{(\lambd[x][Nx])} &= \becd{N} \ollabel{defn:becd5} & \text{که $x
                                                      \notin FV(N)$ وي}
  \end{align}
  د تجريد دپاره که د اېټا شرط پوره وي، وروستۍ قاعده کاروو؛ په نورو تجريدونو کښې عمومي قاعده کاروو۔
\end{defn}

\begin{lem}\ollabel{lem:comp}
  که $M \beredpar M'$ او $R \beredpar R'$ وي، نو $\Subst{M}{R}{y}
  \beredpar \Subst{M'}{R'}{y}$.
\end{lem}

\begin{proof}
  د $M \beredpar M'$ د !!{derivation} پر جوړښت استقرا کوو۔

  په ټول ثبوت کښې تړلي نومونه د الفا-بدلون له لارې تازه ټاکو۔ هممهاله راکمول نوے ازاد نوم نۀ راولي: دا د قاعدو پر جوړښت له استقرا او د تعويض د ازادو نومونو له قاعدې راځي۔ لومړي څلور حالتونه د \olref[pb]{lem:comp} د حالتونو په شان دي۔
  که وروستۍ قاعده \olref{defn:beredpar5} وي، نو $M$،
  $\lambd[x][Nx]$ او $M'$ د ځينو $x$ او~$N'$ دپاره $N'$
  دے، چې $x \notin FV(N)$ او $N \beredpar N'$ وي۔ غواړو
  وښيو چې $\Subst{(\lambd[x][Nx])}{R}{y} \beredpar
  \Subst{N'}{R'}{y}$، يعنې
  $\lambd[x][\Subst{N}{R}{y} x] \beredpar \Subst{N'}{R'}{y}$۔
  دا د \olref{defn:beredpar}، \olref{defn:beredpar5} او
  د استقرا له فرضيې راځي۔
\end{proof}

\begin{lem}\ollabel{lem:cont}
  که $M \beredpar M'$ وي، نو $M' \beredpar \becd{M}$.
\end{lem}

\begin{proof}
  د $M \beredpar M'$ د !!{derivation} وروستۍ قاعده د استقرا په هر ګام کښې څېړو۔

  د دوو دعوو دپاره پر ترمونو يوځايي استقرا کوو، او د کوچني ترم دپاره دواړه دعوې فرضوو۔ لومړۍ دعوه همدا لمه ده۔ دويمه دعوه دا ده: که $P \ident \lambd[x][N]$، $P \beredpar P'$ او $Q' \beredpar \becd{Q}$ وي، نو $P'Q' \beredpar \Subst{\becd{N}}{\becd{Q}}{x}$۔ په دويمه دعوه کښې استقرا د $P$ پر جوړښت ده او د $Q$ او $Q'$ په ټاکلو هېڅ بنديز نشته۔
\par په ټول ثبوت کښې تړلي نومونه د الفا-بدلون له لارې تازه ټاکو۔ هممهاله راکمول نوے ازاد نوم نۀ راولي: دا د قاعدو پر جوړښت له استقرا او د تعويض د ازادو نومونو له قاعدې راځي۔ لومړے دويمه دعوه ثابته کړو۔ که $P'$ د تجريد له قاعدې $\lambd[x][N']$ وي، د لومړۍ دعوې استقرايي فرض له $N \beredpar N'$ څخه $N' \beredpar \becd{N}$ ورکوي۔ اوس د بېټا قاعده پر $P'Q'$ ولګوئ۔ که $P'$ د اېټا له قاعدې راغلے وي، نو $N \ident Sx$، $x \notin FV(S)$ او $S \beredpar P'$۔ که $S$ تجريد نۀ وي، $\becd{N} \ident \becd{S}x$؛ د لومړۍ دعوې استقرايي فرض $P' \beredpar \becd{S}$ ورکوي او د تطبيق قاعده مطلوبه پايله ورکوي۔ که $S \ident \lambd[y][T]$ وي، $\becd{N} \ident \Subst{\becd{T}}{x}{y}$۔ دويمه استقرايي دعوه پر کوچني $S$ ولګوئ: $P'Q' \beredpar \Subst{\becd{T}}{\becd{Q}}{y}$۔ ځکه چې $x$ په $S$ کښې ازاد نۀ دے، دا هدف له $\Subst{\becd{N}}{\becd{Q}}{x}$ سره د تعويض له قاعدې برابر دے۔
\par د لومړۍ دعوې دپاره د متغير حالت انعکاسي دے۔ د تطبيق په حالت کښې که اصلي کيڼ ترم تجريد نۀ وي، د دواړو کوچنيو ترمونو د استقرايي فرضو او د تطبيق له قاعدې هدف راځي۔ که کيڼ ترم تجريد وي او وروستۍ قاعده تطبيق وي، د هغه کيڼ ترم دپاره دويمه استقرايي دعوه کاروو۔ که وروستۍ قاعده د بېټا انقباض وي، د تجريد د بدن او د ښي ترم استقرايي فرضونه د تعويض له لمې سره يوځاے مطلوبه پايله ورکوي۔
\par د تجريد په حالت کښې، که اصلي ترم د اېټا رېډکس نۀ وي، د بدن استقرايي فرض او د تجريد قاعده بس دي۔ که اصلي ترم $\lambd[x][Sx]$ وي او $x \notin FV(S)$، هدف $\becd{S}$ دے۔ که د بدن هممهاله ګام د تطبيق له قاعدې وي، نو پايله $\lambd[x][S'x]$ ده او $S \beredpar S'$۔ له استقرا او د ازادو نومونو له شرطه $S' \beredpar \becd{S}$ او $x \notin FV(S')$ لرو، نو د اېټا قاعده کار کوي۔ که د بدن هممهاله ګام په سر کښې بېټا منقبض کړي، $S \ident \lambd[y][T]$ او $T \beredpar T'$ وي۔ نو پايله $\lambd[x][\Subst{T'}{x}{y}]$ ده، چې د تازګۍ له امله له $\lambd[y][T']$ سره الفا-معادله ده۔ د تجريد قاعده دا له $S$ څخه هممهاله راکموي؛ نو د کوچني $S$ دپاره لومړۍ استقرايي فرضيه دا مستقيماً $\becd{S}$ ته راکموي۔
  که وروستۍ قاعده \olref{defn:beredpar5} وي، نو $M$،
  $\lambd[x][Nx]$ او $M'$ د ځينو $x$، $N$ او $N'$ دپاره
  $N'$ دے، چې $x \notin FV(N)$ او $N \beredpar N'$ وي۔
  غواړو وښيو چې $N' \beredpar
  \becd{(\lambd[x][Nx])}$، يعنې $N' \beredpar \becd{N}$؛
  دا د استقرا له فرضيې سملاسي راځي۔
\end{proof}

\begin{thm}\ollabel{thm:cr}
  $\beredpar$ د چرچ--روسر خاصيت لري۔
\end{thm}

\begin{proof}
  دا د \olref{lem:cont} سملاسي پايله ده۔
\end{proof}

\end{document}
